\begin{afterword}
\summaryhead

The post defines an optimization process as one that hits small targets in a large space of possibilities, and treats natural selection and human intelligence as two such processes. It separates the meta level, the structure that does the optimizing, from the object level, the things produced. Natural selection is simple and began by accident; what it produces, such as a cat, is far more complex. Animal brains were not major optimizers, because their learning neither generalized nor accumulated. Human brains now do both.

In the history of Earth so far, the post argues, the meta level has been protected: genes rarely change natural selection, and science does not rebuild the brain. Rare exceptions such as sex, cells and DNA, and science begin new epochs. An AI that could redesign itself all the way down would have no protected optimizing level, and would break with the whole past. To the objection that each gain might need exponentially more effort, the author answers that roughly constant selection on the hominid line does not seem to have needed exponentially more time for each step of improvement. The author calls all this analogy, but concludes that the right graph is optimization power in against optimized product out, not growth over time.

\responsehead

The post offers a way to describe history, not a proof, and it says so. It calls its argument ``mere analogic reasoning,'' good ``at best'' for qualitative predictions, and ends with a reason ``why I might be reluctant'' to extrapolate growth.

Its central idea has a known source. A machine that improves its own design and so leaves its makers behind is I.~J. Good's ``intelligence explosion'' of 1965. The title uses Good's term, but the post does not mention Good. What the post adds is the account of history as a series of epochs, each with a protected meta level. That account agrees with the biologists' history of major transitions (Maynard Smith and Szathm\'ary), which the post alludes to without naming.

The two premises that connect the account to the conclusion are asserted, not shown. The first is that earlier optimizers worked ``at a constant rate,'' so that acceleration came from innovations opening the way to others and from rare changes to the meta level. The number of optimizers has also grown, and the economist Robin Hanson wrote in 2008 that a larger economy speeds growth by letting more people explore more innovations. The brain stayed the same, but the total optimization put in did not. The second premise is the answer to the objection about diminishing returns. The post's reply is that the hominid line ``doesn't seem'' to have needed exponentially more time for each increment. Neither the constant pressure nor the increment is measured, and the case covers only the climb up to human level. The objection concerns an AI improving itself far beyond that level.

A smaller inference also assumes its conclusion. That rationality training cannot turn an arbitrary person into Einstein is said to show ``the power of a few minor genetic quirks.'' It shows only that training does not close the gap, not that the gap is genetic.

The last paragraph answers a specific forecast that this version does not name. The first version, in June 2008, named Hanson and argued against a forecast of one-month economic doubling times; Hanson had projected past eras of growth forward to a doubling time of a week to a month.

In short: a candid analogy that describes Good's intelligence explosion as the first optimizer able to redesign its own machinery, whose link to its conclusion depends on two unmeasured premises: that past optimization ran at a constant rate, and that the climb from ape to human answers the objection about diminishing returns.
\end{afterword}
